%%
%
%\documentclass[11pt,a4paper,twoside]{article}

\usepackage[utf8]{inputenc}

\usepackage[T1,T2A]{fontenc}

\usepackage[english.ancient,russian.ancient]{babel}

\usepackage{graphicx}

\usepackage{array}

\usepackage[doi]{samgtu-bib}

\usepackage{samgtu}

\usepackage{samgtu-name}

\usepackage{mathtools}

\mathtoolsset{showonlyrefs}

\usepackage{desclist}

\usepackage{euscript}

\usepackage{mathrsfs} 

\usepackage{multicol}

\usepackage{mathrsfs}

\usepackage{cite}

\usepackage{cmap}

\usepackage{qrcode}

\allowdisplaybreaks[4]

\usepackage[nodayofweek]{datetime}

\usepackage[thinlines]{easybmat}



\renewcommand{\JournalYear}{2018}

\renewcommand{\JournalVolume}{22}

\renewcommand{\JournalNumber}{4}



\newdate{JournalDateSign}{29}{12}{2018}% Дата подписания Вестника в печать

\renewcommand{\JournalDateSign}{\displaydate{JournalDateSign}}

\newdate{JournalDatePrint}{16}{01}{2019}% Дата выхода Вестника из печати

\renewcommand{\JournalDatePrint}{\displaydate{JournalDatePrint}}

%\renewcommand{\JournalUPL}{16.56} % Усл. печ. л.

%\renewcommand{\JournalUIL}{16.53} % Уч.-изд. л.

\renewcommand{\JournalUPL}{15.24} % Усл. печ. л.

\renewcommand{\JournalUIL}{15.21} % Уч.-изд. л.

\renewcommand{\JournalRegNumber}{220/18} % Регистрационный номер

\renewcommand{\JournalOrderNumber}{2} % Номер заказа



%\renewcommand{\JournalRMonth}{Март} % Месяц выхода Вестника

%\renewcommand{\JournalEMonth}{March} % Месяц выхода Вестника

%\renewcommand{\JournalRMonth}{Июнь} % Месяц выхода Вестника

%\renewcommand{\JournalEMonth}{June} % Месяц выхода Вестника

%\renewcommand{\JournalRMonth}{Сентябрь} % Месяц выхода Вестника

%\renewcommand{\JournalEMonth}{September} % Месяц выхода Вестника

\renewcommand{\JournalRMonth}{Декабрь} % Месяц выхода Вестника

\renewcommand{\JournalEMonth}{December} % Месяц выхода Вестника



\newcommand{\totop}[1]{\raisebox{6pt}[1pt][2pt]{#1}}

\newcommand{\tottop}[1]{\raisebox{12pt}[1pt][2pt]{#1}} 

\newcommand{\totttop}[1]{\raisebox{18pt}[1pt][2pt]{#1}} 

\newcommand{\tottttop}[1]{\raisebox{24pt}[1pt][2pt]{#1}} 

\newcommand{\totttttop}[1]{\raisebox{30pt}[1pt][2pt]{#1}} 

\newcommand{\tobot}[1]{\raisebox{-6pt}[1pt][2pt]{#1}} 

\newcommand{\tobbot}[1]{\raisebox{-12pt}[1pt][2pt]{#1}} 

\newcommand{\tobbbot}[1]{\raisebox{-18pt}[1pt][2pt]{#1}}



\graphicspath{{./00PIC/}}
%
%\usepackage{samgtu-name}
%\usepackage{wrapfig}
%
%
\vsgtu{1613}
%
%
\Personalia
%\pdfcrop % обрезка pdf для размещения на math-net.ru
%\draft % На корректуре
%
%
%
\FirstPage{15}
\LastPage{22}
%

%\makeatletter
%\newcommand{\PersonRef}{\@langrustrue\def\refname{{\bf\small Избранные  труды А.~П.~Солдатова за последние 5 лет}\footnote{Список подготовлен М.~Н.~Саушкиным (Самара, СамГТУ).}}}
%%
%\newcommand{\ManzhirovRef}{\@langrustrue\def\refname{{\bf\small Статьи о проф. А.~В.~Манжирове}}}
%\newcommand{\ManzhirovRefEng}{\@langrustrue\def\refname{{\bf\small Publications  about Prof. A.~V.~Manzhirov in the Press}}}
%
%\makeatother



%\begin{document}
$\;$

\hfill \qrcode[hyperlink,height=.8in]{http://doi.org/10.14498/vsgtu1613}
\vspace{-.8in}



\rutitle{К 70-летию профессора Александра Павловича Солдатова}
\rutitleshort{К 70-летию профессора Александра Павловича Солдатова}
\rutitlehack{К 70-летию профессора \\Александра Павловича Солдатова}
\entitle{To the 70$^{\rm th}$ Anniversary of Professor Alexander Pavlovich Soldatov}
\entitlehack{To the 70$^{\bf th}$ Anniversary of\\ Professor Alexander Pavlovich Soldatov}

\ruauthor{А.~А.~Андреев, В.~П.~Радченко, Е.~А.~Козлова}{А\,н\,д\,р\,е\,е\,в~А.~А., Р\,а\,д\,ч\,е\,н\,к\,о~В.~П., К\,о\,з\,л\,о\,в\,а~Е.~А.}
\enauthor{A.~A.~Andreyev, V.~P.~Padchenko, E.~A.~Kozlova}{A\,n\,d\,r\,e\,y\,e\,v~A.~A., R\,a\,d\,c\,h\,e\,n\,k\,o~V.~P., K\,o\,z\,l\,o\,v\,a~E.~A.}

\ruaffil{\rusamgtu.}
\enaffil{\ensamgtu.}

\selectlanguage{russian}%
\phantomsection%
\addtocontent{{\it\RuAuthorInContents}\hspace{2.5mm}}{``\RuTitleInContents''}
\addtoengcontent{{\it\EnAuthorInContents}\hspace{2.5mm}}{``\selectlanguage{english}\EnTitleInContents''}
\normalsize 
\normalfont
\chead[\fancyplain{}{\selectlanguage{russian}\theruauthor}]{\fancyplain{}{\selectlanguage{russian} \therutitle}}
\rhead[]{\fancyplain{\RuJournal}{}}
\lhead[\fancyplain{\RuJournal}{}]{}
\cfoot[]{}
\lfoot[\thepage]{}
\rfoot[]{\thepage}
\renewcommand{\plainheadrulewidth}{0.7pt}
\renewcommand{\headrulewidth}{0.7pt}
\pagestyle{fancyplain}
\thispagestyle{plain}
\clear
\hskip-7mm
\begin{minipage}[t]{125mm}
\vskip.5mm
\noindent\theudc
\vskip2mm
\RuTitleInTitlepage
\end{minipage}
\par
\theruauthorsdetails
\begin{minipage}[t]{125mm}
\vskip4mm
{\renewcommand{\bfdefault}{bx}\noindent\small\textbf{\emph{\RuAuthorInTitlepage}}\renewcommand{\bfdefault}{b}}  
\vskip2mm
{\scriptsize \ruaffilname}
\end{minipage}
\vskip2mm


\footnotetext{
\vspace{-1em}\\
\noindent{\normalsize\bf\ruArticleType}\smallskip\par\noindent
\OAlogo~\ccby\  Контент публикуется на условиях  лицензии  \href{https://creativecommons.org/licenses/by/4.0/deed.ru}{Creative Commons Attribution 4.0 International} (\url{https://creativecommons.org/licenses/by/4.0/deed.ru})
\vspace{.3em}\\ 
\RuCitation}


\begin{wrapfigure}{o}{0.46\textwidth} 
\vspace{-3mm}
\centering
\includegraphics[width=0.4\textwidth]{soldatov2.jpg}

\vspace{-8mm}

\end{wrapfigure} 

\footnotesize

\smallskip


18 января 2018~г. исполнилось 70 лет доктору физико-математических наук, профессору, заслуженному деятелю науки Российской Федерации Александру Павловичу Солдатову.

В настоящей  заметке приводятся сведения из научной биографии этого крупного ученого, который хорошо известен в~России и за рубежом. 
Представлены сведения о~его вкладе в~развитие математики в области дифференциальных уравнений с~частными производными и~их приложений.


\newdate{soldatove}{29}{03}{2018}
\renewcommand{\JournalDataOnline}{\displaydate{soldatove}}

\smallskip

\par\noindent
\theonlinerudate \par\noindent
\rule{68mm}{.7pt}\\

\normalsize


Александр Павлович Солдатов родился 18 января 1948 года в~Тюменской области. 
Среднее образование он получил в школе"=интернате при \ruaffid{2794}{Новосибирском государственном университете}, в~котором и~продолжил дальнейшее образование, окончив в~1971 году с~отличием университет по специальности <<Математика>>. 
В~1975 году защитил кандидатскую диссертацию в~\ruaffid{748}{Математическом институте АН СССР им. В.~А.~Стеклова}, а~в~1981 году "--- докторскую диссертацию в~\ruaffid{189}{Московском государственном университете им. М.~В.~Ломоносова}. 
Имеет ученое звание профессора с 1985 года, заслуженный деятель науки Российской Федерации, действительный член Международной академии наук высшей школы с~2001 года. 

А.~П.~Солдатов является известным и активно работающим специалистом в области дифференциальных уравнений с~частными производными, автором более 170 научных статей и~четырех монографий.

Первый цикл работ А.~П.~Солдатова посвящен уравнениям смешанного эллиптико"=гиперболического типа. 
Им впервые была установлена однозначная разрешимость задачи Франкля и~обобщенная задача Трикоми для уравнения Лаврентьева"--~Бицадзе без каких"=либо ограничений геометрического характера на эллиптическую часть границы смешанной области. 
Данная проблема, поставленная \rupid{20250}{А.~В.~Бицадзе}, оставалась открытой начиная с~50-х годов прошлого столетия. 
Актуальность этой тематики в первую очередь определилась потребностями аэродинамики больших скоростей. 
Пионерские работы в области трансзвуковой динамики, приводящие к~новым постановкам краевых задач, по праву принадлежат \rupid{27239}{Ф.~И.~Франклю} и {Г.~Гудерлею} (\href{https://zbmath.org/authors/?q=ai:guderley.gottfried}{Gottfried Guderley}). 

Сравнительно недавно (2009--2010 гг.) А.~П.~Солдатовым предложены новые корректные постановки смешанных задач для модельного уравнения Лаврентьева"--~Бицадзе. 
Эти постановки тесно связаны с~так называемой <<околозвуковой>> полемикой, которая развернулась в середине 50-х между математиками и механиками (Г.~Гудерлей, \rupid{26652}{Г.~Буземан}, Ф.~И.~Франкль и~др.) по вопросу о~существовании гладких околозвуковых решений задачи обтекания произвольно заданного профиля. 
Положительное решение этого вопроса тесно связано с~разрешимостью задачи Дирихле во всей смешанной области. 
Как показано А.~П.~Солдатовым, эта задача однозначно разрешима в~классе классических решений, допускающих особенности сколь угодно малого порядка в~одной из граничных точек на линии изменения типа. 
Следует отметить, что в~классе обобщенных решений близкие результаты были независимо получены известным американским математиком {К.~Моравец} (\href{https://zbmath.org/authors/?q=ai:synge-morawetz.cathleen}{Cathleen Morawetz}).

Значительное число работ А.~П.~Солдатова связано с разработкой аппарата одномерных сингулярных интегро"=функциональных уравнений. 
Им выделен важный для приложений класс уравнений этого типа и построена фредгольмовская теория операторов данного класса; найдены критерий фредгольмовости и~формула индекса, а~также изучены вопросы асимптотики и~гладкости решений уравнений в окрестности особых точек линии интегрирования. 
Итог этим исследованиям подведен в~монографии:
\href{https://zbmath.org/?q=an:0774.47025}{
{\it Солдатов~А.~П.} Одномерные сингулярные операторы и краевые задачи теории функций. М.: Высшая школа, 1991. 208~с.} 

К данному   направлению   примыкает третий цикл работ А.~П.~Солдатова, в~котором предложен и~развит новый подход к~исследованию общих (локальных и нелокальных) эллиптических краевых задач на плоскости. 
Он заключается в~эквивалентной редукции краевой задачи к~сингулярному интегро"=функциональному уравнению на граничном контуре области, основанной на представлении решений эллиптических уравнений и~систем через так называемые функции, аналитические по Дуглису. 
Посредством этой редукции для эллиптических задач получаются результаты той же степени полноты, что и~для интегральных уравнений на границе. 
Этот метод нашел свое отражение в~двухтомной монографии:
\href{https://zbmath.org/?q=an:0854.30029}{{\it Солдатов~А.~П.} Краевые задачи теории функций в областях с кусочно-гладкой границей. Тбилиси: Тбилисский государственный университет, 1991. Т.~1. 268~с.; Т.~2. 276~с.}


Ряд важных результатов получен им в~плоской теории упругости, где предложены новые интегральные представления общего решения эллиптической системы Ламе, которые можно рассматривать как варианты потенциалов двойного слоя. 
Эти представления позволяют осуществить новый способ редукции основных краевых задач к~эквивалентной системе интегральных уравнений Фредгольма на границе области. 
Аналогичный подход оказалось возможным развить и~для общих эллиптических систем второго порядка, для которых задача Дирихле фредгольмова. 
Еще в~80-х годах А.~П.~Солдатовым была исследована фредгольмова разрешимость так называемой общей нелокальной задачи Римана для аналитических функций, обобщающая классическую задачу Римана"--~Гильберта. 
В~последнее время неожиданное применение она нашла в~решении задачи Дирихле для гармонических функций в~двумерных стратифицированных областях. 
Эти результаты получены в~последнее десятилетие и~опубликованы в~нескольких центральных математических журналах. 

Среди совсем недавних результатов отметим полученное А.~П.~Солдатовым решение в явном виде обратной задачи для уравнения Штурма"--~Лиувилля на всей оси. 
Общая проблема существования и единственности этой задачи была исследована в~классических работах 
\rupid{8938}{И.~М.~Гельфанда}, 
\rupid{20095}{Б.~М.~Левитана}, 
\rupid{9160}{Л.~Д.~Фаддеева}, 
\rupid{15243}{В.~А.~Марченко} и~основывается на использовании интегральных уравнений, носящих их имя. 
В~совместной 
работе\footnote{{\textit{Жура~Н.~А., Солдатов~А.~П.} К решению обратной задачи Штурма--Лиувилля на всей оси /\!/ Доклады Акад. наук. 2013. Т.~453, № 4. С.~368--372.} doi: \href{http://dx.doi.org/10.7868/S0869565213340070}{10.7868/S0869565213340070}}\! А.~П.~Солдатова и~\rupid{20611}{Н.~А.~Жура}, во-первых, указаны достаточные условия выполнимости основной теоремы Фаддеева"--~Марченко, о которой шла речь выше, и,~во-вторых, дано представление решения обратной задачи на основе решения краевой задачи для функций Йоста и отвечающего ей сингулярного интегрального уравнения. Последний подход не нов и рассматривался, например, в работах 
\rupid{18696}{В.~Е.~Захарова} и~\rupid{18749}{А.~Б.~Шабата}, где и было выведено упомянутое сингулярное уравнение. 
Однако это уравнение не было ими исследовано. 

Несколько в стороне от основного направления исследований лежит предложенная А.~П.~Солдатовым (совместно с Н.~А.~Жура) постановка корректной краевой задачи для гиперболических систем на 
плоскости.\!\footnote{{\textit{Жура~Н.~А., Солдатов~А.~П.} Краевая задача для гиперболической системы первого порядка в двумерной области/\!/ Изв. РАН. Сер. матем. 2017. Т.~81, №~3. C.~83–108.} doi: \href{http://doi.org/10.1070/im8442}{10.1070/im8442}.}\!\!~Интерес к~краевым задачам для гиперболических уравнений возник в~начале тридцатых годов прошлого столетия, хотя задачи типа задачи Гурса, которые часто возникают в~приложениях, были известны и~ранее. 
Различные ее обобщения для систем второго порядка гиперболического типа в~бесконечных областях типа сектора исследовались А.~В.~Бицадзе, \rupid{31179}{С.~С.~Харибегашвили} и рядом других авторов. Однако в конечных областях систематическому исследованию было подвергнуто лишь уравнение для струны. 

Исследования А.~П.~Солдатова привлекли широкое внимание специалистов как у нас в стране, так и за рубежом. Результаты исследований только в последние несколько лет докладывались и обсуждались на различных международных симпозиумах и конгрессах:
\begin{itemize}
\item International Conference on {\it Differencial equations}, December 9--12, 2017, Dalat, 
Vietnam;

\item[--] International Conference ``{\it Applications of Mathematics in Engineering and Economics}'' (AMEE'17), June 8--13,  2017, Sozopul, Bulgaria; 
\item[--] International Conference on {\it Analysis and Applied Mathematics} (ICAAM 2016), September 7--10, 2016, Almaty, Kazakhstan; 
\item[--] 
3rd International Workshop  ``{\it Boundary value problems, functional equations and their applications}'', April 20--23, 2016, Rzesz\'ow,  Poland;
\item[--] International Workshop ``{\it Wiener--Hopf Method, Toeplitz Operators, and their Applications}'',
November 3--7, 2015,  Veracruz, Mexico;
\item[--] International Conference on {\it Mathematical Problems in Engineering, Aero\-space and Sciences} (ICNPAA 2014), July 15–18, 2014, Narvik, Norway;

\item[--] International Conference on {\it Finite or Infinite Dimensional Complex Analysis and Applications},
June 15--19, 2013, Nanjing, China;
\item[--] International Scientific Conference ``{\it Continuum Mechanics and Analysis of the Related Problems}'', September 10--14, 2011, Tbilisi, Georgia и многих других.

\end{itemize}

Профессор А.~П.~Солдатов является членом ряда отечественных и международных математических обществ и редколлегий научных журналов. Он член  \href{http://mms.mathnet.ru/}{Московского математического общества},
Американского математического общества (\href{http://www.ams.org/home/page}{AMS}), 
Международной ассоциации по прикладной математике и механике (\href{https://www.gamm-ev.de/}{GAMM}), 

%Он ведет активную организаторскую и научно"=педагогическую работу. 
%Во время работы на факультете математики и~информационных технологий  \ruaffid{2866}{Белгородского государственного университета} на основе преподаваемого в~течение ряда лет университетского 4-х семестрового курса по математическому анализу Александром Павловичем подготовлен новый учебник по этой базовой дисциплине, который вышел отдельными частями в~виде учебных пособий. Концепция построения этого курса, учитывающая современное состояние науки, неоднократно обсуждалась на различных крупных конференциях научно"=методического профиля. 

Он ведет активную организаторскую и научно"=педагогическую работу, принимает активное участие в~различных программах в~области  школьного  математического  образования.  В~течение длительного  времени  Александр Павлович являлся  председателем  жюри  математических олимпиад Владимирской области и~был отмечен почетным знаком <<Отличник просвещения>>.\!\footnote{Решение № 99-к от 10.07.90 г. Министерства просвещения РСФСР.}

\vspace{-3mm}

По инициативе А.~П.~Солдатова в~2006 году был открыт диссертационный совет по защите диссертаций на соискание ученой степени доктора и~кандидата наук по специальности <<Дифференциальные уравнения>>, председателем которого он является. 
Этот совет играет заметную роль в развитии научных исследований  и~подготовке кадров высшей квалификации. 
С~его деятельностью тесно связан научный семинар по дифференциальным уравнениям и~краевым задачам, на котором проходят экспертизу многие кандидатские и~докторские диссертации. 
А.~П.~Солдатов принимает активное участие в консультировании, рецензировании и~оппонировании кандидатских и~докторских диссертаций. 
Только за последние 5 лет он являлся официальным оппонентом свыше  десяти  докторских и кандидатских диссертаций.

А.~П.~Солдатов в течение двух десятков лет тесно сотрудничает с~\ruaffid{2786}{Самарским государственным техническим университетом}, являясь членом редколлегии журнала <<\href{http://doi.org/10.14498/vsgtu}{Вестник Самарского государственного технического университета. Сер. Физико-ма\-те\-ма\-ти\-чес\-кие науки}>>, членом оргкомитетов более десятка всероссийских и международных конференций в области дифференциальных уравнений, уравнений в частных производных и математического моделирования, которые проходили на базе технического университета.
Многие преподаватели и научные сотрудники технического университета в свое время защитили свои диссертации на соискание ученой степени кандидата наук в диссертационном совете, председателем которого является А.~П.~Солдатов.

Профессор А. П. Солдатов "--- ученый высочайшей квалификации в области уравнений частных производных и~их приложений, о чем свидетельствуют его научные работы, опубликованные за последние несколько лет, которые приводятся ниже. 
Его отличает требовательное, но доброжелательное отношение к ученикам и коллегам.
Он всегда занимает принципиальную позицию в научных дискуссиях и во главу ставит прежде всего вопросы поиска научной истины.




\newpage
\PersonRef

\begin{thebibliography}{99}

\RBibitem{1}
\by Солдатов~А.~П., Кожанов~А.~И.
\paper \href{http://mi.mathnet.ru/vngu441}{А. В. Бицадзе. К столетию со дня рождения}
\jour Сиб. журн. чист. и прикл. матем.
\yr 2017
\vol 17
\issue 3
\pages 3--7
\mathnet{http://mi.mathnet.ru/vngu441}

\RBibitem{1}
\by Жура~Н.~А., Солдатов~А.~П.
\paper Краевая задача для гиперболической системы первого порядка в двумерной области
\jour Изв. РАН. Сер. матем.
\yr 2017
\vol 81
\issue 3
\pages 83--108
\mathnet{http://mi.mathnet.ru/izv8442}
\crossref{10.4213/im8442}
\mathscinet{http://www.ams.org/mathscinet-getitem?mr=3659547}
\adsnasa{http://adsabs.harvard.edu/cgi-bin/bib_query?2017IzMat..81..542Z}
\elib{http://elibrary.ru/item.asp?id=29254883}
\transl
\paper A boundary-value problem for a first-order hyperbolic system in a two-dimensional domain
\by Zhura~N.~A.,  Soldatov~A.~P.
\jour Izv. Math.
\yr 2017
\vol 81
\issue 3
\pages 542--567
\crossref{10.1070/IM8442}
\isi{http://gateway.isiknowledge.com/gateway/Gateway.cgi?GWVersion=2&SrcApp=PARTNER_APP&SrcAuth=LinksAMR&DestLinkType=FullRecord&DestApp=ALL_WOS&KeyUT=000408479100004}
\scopus{http://www.scopus.com/record/display.url?origin=inward&eid=2-s2.0-85025477852}


\RBibitem{1}
\by Полунин~В.~А., Солдатов~А.~П.
\paper Об интегральном представлении решений системы Моисила--Тедореску в многосвязных областях
\jour Докл. Акад. наук
\crossref{10.7868/S0869565217220029}
\yr 2017
\vol 475
\issue 4
\pages 369–372
\transl
\by Polunin~V.~A.,  Soldatov~A.~P.
\jour Dokl. Math. 
\crossref{10.1134/S1064562417040172}
\yr 2017
\vol 96
\issue 1
\pages 358–361 
\paper Integral representation of solutions of the Moisil–Théodorescu system in multiply connected domains


\RBibitem{Sol17}
\by Солдатов~А.~П.
\paper Сингулярные интегральные операторы и эллиптические краевые задачи.~I
\inbook Функциональный анализ
\serial СМФН
\yr 2017
\vol 63
\issue 1
\pages 1--189
\publ РУДН
\publaddr М.
\mathnet{http://mi.mathnet.ru/cmfd316}
\crossref{10.22363/2413-3639-2017-63-1-1-189}
\mathscinet{http://www.ams.org/mathscinet-getitem?mr=3646636}

\RBibitem{1}
\by Солдатов~А.~П.
\paper \href{http://mi.mathnet.ru/vmj624}{Об одной краевой задаче для эллиптического уравнения высокого порядка в многосвязной области на плоскости}
\jour Владикавк. матем. журн.
\yr 2017
\vol 19
\issue 3
\pages 51--58
\mathnet{http://mi.mathnet.ru/vmj624}

\RBibitem{1}
\crossref{10.1134/S0374064116090028}
\paper Асимптотика решений задачи линейного сопряжения для аналитических функций в угловых точках кривой
\by Аверьянов~Г.~Н., Солдатов~А.~П.
\jour Диффер. уравн.
\yr 2016
\vol 52
\issue 9
\pages 1150--1159
\transl
\by Aver'yanov~G.~N.,  Soldatov~A.~P.
\paper Asymptotics of solutions of the linear conjugation problem at the corner points of the curve
\jour Differ. Equat.
\vol 52
\issue 9
\pages 1105--1114 
\yr 2016
\crossref{10.1134/S0012266116090019}

\RBibitem{1}
\crossref{10.1134/S0374064116040105}
\by Мещерякова~Е.~С., Солдатов~А.~П.
\paper Задача Римана--Гильберта в семействе весовых пространств Гельдера
\jour Диффер. уравн.
\yr 2016
\vol 52
\issue 1
\pages 518--527 
\transl 
\paper Riemann–Hilbert problem in a family of weighted Hölder spaces
\by Meshcheryakova~E.~S., Soldatov~A.~P.
\jour Differ. Equat.
\yr 2016
\vol 52
\issue 4
\pages 495--504
\crossref{10.1134/S0012266116040091}

\RBibitem{1}
\crossref{10.1134/S0374064116120074}
\by Кошанов~Б.~Д, Солдатов~А.~П.
\paper Краевая задача с нормальными производными для эллиптического уравнения на плоскости
\jour Диффер. уравн.
\yr 2016
\vol 52
\issue 12
\pages 1666--1681
\transl
\paper Boundary value problem with normal derivatives for a higher-order elliptic equation on the plane
\by Koshanov~B.~D., Soldatov~A.~P.
\jour Differ. Equat.
\yr 2016
\vol 52
\issue 12
\pages 1594--1609
\crossref{10.1134/S0012266116120077}

\RBibitem{1}
\by Расулов~А.~Б., Солдатов~А.~П.
\crossref{10.1134/S0374064116050083}
\paper Краевая задача для обобщенного уравнения Коши--Римана с~сингулярными коэффициентами
\jour Диффер. уравн.
\yr 2016
\vol 52
\issue 5
\pages 637--650
\transl
\paper Boundary value problem for a generalized Cauchy–Riemann equation with singular coefficients
\by Rasulov~A.~B., Soldatov~A.~P.
\jour Differ. Equat.
\yr 2016
\vol 52
\issue 5
\pages 616--629
\crossref{10.1134/S0012266116050086}

\Bibitem{1}
\by Polunin~V.~A., Soldatov~A.~P.
\paper \href{https://ejde.math.txstate.edu/Volumes/2016/310/polunin.pdf}{Riemann--Hilbert problem for the Moisil--Teodorescu system in multiply connected domains}
\jour Electronic journal of differential equation
\papernumber  310
\yr 2016
\vol 2016
\totalpages 5


\RBibitem{Sol16}
\by Солдатов~А.~П.
\paper О спектральном радиусе функциональных операторов
\jour Матем. заметки
\yr 2016
\vol 100
\issue 1
\pages 155--162
\mathnet{http://mi.mathnet.ru/mz11128}
\crossref{10.4213/mzm11128}
\mathscinet{http://www.ams.org/mathscinet-getitem?mr=3588835}
\elib{http://elibrary.ru/item.asp?id=26414286}
\transl
\paper On the spectral radius of functional operators
\by Soldatov~A.~P.
\jour Math. Notes
\yr 2016
\vol 100
\issue 1
\pages 132--138
\crossref{10.1134/S0001434616070129}
\isi{http://gateway.isiknowledge.com/gateway/Gateway.cgi?GWVersion=2&SrcApp=PARTNER_APP&SrcAuth=LinksAMR&DestLinkType=FullRecord&DestApp=ALL_WOS&KeyUT=000382193300012}
\scopus{http://www.scopus.com/record/display.url?origin=inward&eid=2-s2.0-84983795590}

\RBibitem{16}
\by Солдатов~А.~П.
\crossref{10.1134/S0374064116060108}
\paper Смешанная задача плоской ортотропной теории упругости в полуплоскости
\jour Диффер. уравн.
\yr  2016
\vol 52
\issue 6
\pages 820--833
\transl
\paper Mixed problem of plane orthotropic elasticity in a half-plane
\by Soldatov~A.~P.
\jour Differ. Equat.
\yr  2016
\vol 52
\issue 6
\pages 798--812
\crossref{10.1134/S0012266116060100}


\RBibitem{Sol16}
\by Солдатов~А.~П.
\paper \href{http://mi.mathnet.ru/cmfd298}{К теории анизотропной плоской упругости}
\inbook Труды Седьмой Международной конференции по дифференциальным и функционально-дифференциальным уравнениям \rm (Москва, 22--29 августа, 2014). Часть~3
\serial СМФН
\yr 2016
\vol 60
\pages 114--163
\publ РУДН
\publaddr М.
\mathnet{http://mi.mathnet.ru/cmfd298}



\RBibitem{Sol16}
\by Солдатов~А.~П., Выонг~Чан~К.
\paper \href{http://mi.mathnet.ru/ivm9188}{Задача линейного сопряжения для бианалитических функций}
\jour Изв. вузов. Матем.
\yr 2016
\issue 12
\pages 76--81
\mathnet{http://mi.mathnet.ru/ivm9188}
\transl 
\paper The linear conjugation problem for bi-analytic functions
\by  Soldatov~A.~P., Vuong~Tran Quang
\jour Russian Math. (Iz. VUZ)
\yr 2016
\vol 60
\issue 12
\pages 62--66
\crossref{10.3103/S1066369X16120094}
\isi{http://gateway.isiknowledge.com/gateway/Gateway.cgi?GWVersion=2&SrcApp=PARTNER_APP&SrcAuth=LinksAMR&DestLinkType=FullRecord&DestApp=ALL_WOS&KeyUT=000409309100009}
\scopus{http://www.scopus.com/record/display.url?origin=inward&eid=2-s2.0-84997294287}


\Bibitem{1}
\crossref{10.4314/jfas.v8i2s.626}
\by Polunin~V.~A., Soldatov~A.~P.
\paper An analogue of the Schwarz problem for the Moisil--Teodorescu system in a multiply connected domain
\jour J. Fundam. Appl. Sci.
\yr  2016
\vol 8 
\issue 2S
\pages 2989--2995

\RBibitem{1}
\by Солдатов~А.~П.
\paper \href{http://elibrary.ru/item.asp?id=29201715}{А.~В.~Бицадзе. Достойное служение науке (к столетию со дня рождения)}
\jour Научные ведомости БелГУ. Математика. Физика
\yr 2016
\vol 45
\issue 27
\pages 5--9


\RBibitem{1}
\by Полунин~В.~А., Солдатов~А.~П.
\paper \href{http://elibrary.ru/item.asp?id=29201716}{Система Моисила–Теодореску в многосвязных областях}
\jour Научные ведомости БелГУ. Математика. Физика
\yr 2016
\vol 45
\issue 27
\pages 10--15


\RBibitem{1}
\paper \href{https://elibrary.ru/item.asp?id=26723918}{О граничных свойствах конформных отображений областей с кусочно-гладкой границей}
\by Мещерякова~Е.~С., Солдатов~А.~П.
\jour Известия Юго-Западного государственного университета. Сер. Техника и технологии
\yr 2016
\issue 2(19)
\pages 121--127

\RBibitem{1}
\paper \href{https://elibrary.ru/item.asp?id=25942190}{Интегралы типа потенциалов в весовых пространствах на плоскости}
\by Солдатов~А.~П., Чернова~О.~В.
\jour Известия Юго-Западного государственного университета. Сер. Техника и технологии
\yr 2016
\issue 1(18)
\pages  100--103

\RBibitem{1}
\crossref{10.7868/S086956521513006X}
\by Солдатов~А.~П.
\paper Обобщенные потенциалы двойного слоя анизотропной плоской теории упругости
\jour Докл. Акад. наук
\yr 2015
\vol 462
\issue 1
\pages 526--529
\transl
\crossref{10.1134/S1064562415030047}
\paper Generalized double-layer potentials in anisotropic elasticity on the plane
\by Soldatov~A.~P.
\jour Dokl. Math.
\yr  2015
\vol 91
\issue 3
\pages 269--272

\RBibitem{1}
\by Солдатов~А.~П.
\crossref{10.7868/S0869565215230085}
\paper  Явное описание обобщенных потенциалов двойного слоя для системы Ламе
\jour Докл. Акад. наук
\yr 2015
\vol 463
\issue 5
\pages 21--225
\transl
\crossref{10.1134/S1064562415040262}
\paper Explicit description of generalized double-layer potentials for the Lamé system
\by \linebreak Soldatov~A.~P.
\jour Dokl. Math.
\yr  2015
\vol  92
\issue 1
\pages 487--490

\Bibitem{1}
\by Soldatov~A.~P.
\paper A plane elasticity analogue of the Keldish--Sedov formula
\jour AIP Conf. Proc.
\crossref{10.1063/1.4936711}
\vol 1690
\papernumber 040004 
\yr 2015

\RBibitem{1}
\by Николаев~В.~Г., Солдатов~А.~П.
\crossref{10.1134/S0374064115070158}
\paper О решении задачи Шварца для $J$-аналитических функций в областях, ограниченных контуром Ляпунова
\jour Диффер. уравн.
\yr 2015
\vol 51
\issue 7
\pages 965--969
\transl
\paper On the solution of the Schwarz problem for $J$-analytic functions in a domain bounded by a Lyapunov contour
\by Nikolaev~V.~G., Soldatov~A.~P.
\jour Differ. Equat.
\yr  2015
\vol 51
\issue 7
\pages 962--966
\crossref{10.1134/S0012266115070150}

\RBibitem{KovSol15}
\by Ковалева~Л.~А., Солдатов~А.~П.
\paper Задача Дирихле на двумерных стратифицированных множествах
\jour Изв. РАН. Сер. матем.
\yr 2015
\vol 79
\issue 1
\pages 77--114
\mathnet{izv8223}
\crossref{10.4213/im8223}
\mathscinet{http://www.ams.org/mathscinet-getitem?mr=3352583}
\zmath{https://zbmath.org/?q=an:06428106}
\adsnasa{http://adsabs.harvard.edu/cgi-bin/bib_query?2015IzMat..79...74K}
\elib{http://elibrary.ru/item.asp?id=23421415}
\transl
\paper The Dirichlet problem on two-dimensional stratified sets
\by \linebreak Kovaleva~L.~A., Soldatov~A.~P.
\jour Izv. Math.
\yr 2015
\vol 79
\issue 1
\pages 74--108
\crossref{10.1070/IM2015v079n01ABEH002735}
\isi{http://gateway.isiknowledge.com/gateway/Gateway.cgi?GWVersion=2&SrcApp=PARTNER_APP&SrcAuth=LinksAMR&DestLinkType=FullRecord&DestApp=ALL_WOS&KeyUT=000350754500005}
\scopus{http://www.scopus.com/record/display.url?origin=inward&eid=2-s2.0-84924284439}

\RBibitem{KovSol15}
\crossref{10.1134/S0374064115080063}
\by Жура~Н.~А., Солдатов~А.~П.
\paper О представлении решения обратной задачи Штурма–Лиувилля на всей оси
\jour Диффер. уравн.
\yr 2015
\vol 51
\issue 8
\pages 1027--1037 
\transl
\paper On a representation of the solution of the inverse Sturm–Liouville problem on the entire line
\by Zhura~N.~A., Soldatov~A.~P.
\jour Differ. Equat.
\yr 2015
\vol 51
\issue 8
\pages 1022--1032
\crossref{10.1134/S0012266115080066}

\RBibitem{AveSol15}
\by Аверьянов~Г.~Н., Солдатов~А.~П.
\paper \href{http://mi.mathnet.ru/ivm9035}{Задача линейного сопряжения для аналитических функций в~семействе весовых пространств Гёльдера}
\jour Изв. вузов. Матем.
\yr 2015
\issue 9
\pages 56--61
\mathnet{http://mi.mathnet.ru/ivm9035}
\transl
\paper Linear conjugation problem for analytic functions in the weighted Hölder spaces
\by Aver'yanov~G.~N., Soldatov~A.~P.
\jour Russian Math. (Iz. VUZ)
\yr 2015
\vol 59
\issue 9
\pages 47--50
\crossref{10.3103/S1066369X15090066}
\scopus{http://www.scopus.com/record/display.url?origin=inward&eid=2-s2.0-84943257511}

\Bibitem{Sol13}
\by Zhura~N.~A., Soldatov~A.~P.
\paper A boundary value problem for first order strictly hyperbolic systems on the plane
\jour AIP Conf. Proc.
\vol 1690
\papernumber 040015 
\yr 2015
\crossref{10.1063/1.4936722}

\RBibitem{Sol13}
\by Аверьянов~Г.~Н., Солдатов~А.~П.
\paper \href{https://elibrary.ru/item.asp?id=24194784}{Асимптотика решений задачи линейного сопряжения для аналитических функций в угловых точках кривой}
\jour Научные ведомости БелГУ. Математика. Физика
\yr 2015
\vol 38
\issue 5
\pages 5--17

\RBibitem{Sol13}
\by Солдатов~А.~П., Выонг~Чан~К. 
\paper \href{https://elibrary.ru/item.asp?id=24194792}{Задача Римана--Гильберта для бианалитических функций}
\jour Научные ведомости БелГУ. Математика. Физика
\yr 2015
\vol 38
\issue 5
\pages 83–88 

\Bibitem{Sol14}
\by A.~P.~Soldatov
\paper \href{http://mi.mathnet.ru/emj158}{Generalized potentials of double layer in plane theory of elasticity}
\jour Eurasian Math.~J.
\yr 2014
\vol 5
\issue 2
\pages 78--125
\mathnet{http://mi.mathnet.ru/emj158}

\RBibitem{Sol13}
\by Солдатов~А.~П.
\paper \href{http://mi.mathnet.ru/cmfd244}{Задача Неймана для эллиптических систем на плоскости}
\inbook Труды Шестой Международной конференции по дифференциальным и функционально-дифференциальным уравнениям \rm (Москва, 14--21 августа, 2011). Часть~4
\serial СМФН
\yr 2013
\vol 48
\pages 120--133
\publ РУДН
\publaddr М.
\mathnet{http://mi.mathnet.ru/cmfd244}
\transl
\paper The Neumann Problem for Elliptic Systems on a Plane
\by Soldatov~A.
\jour J. Math. Sci.
\yr 2014
\vol 202
\issue 6
\pages 897--910
\crossref{10.1007/s10958-014-2085-7}
\scopus{http://www.scopus.com/record/display.url?origin=inward&eid=2-s2.0-84919916313}


\RBibitem{Sol13}
\by Митин~С.~П., Солдатов~А.~П.
\paper \href{https://elibrary.ru/item.asp?id=21639537}{Задача Дирихле для системы Ламе с кусочно-пос\-тоян\-ны\-ми коэффициентами}
\jour Научные ведомости БелГУ. Математика. Физика
\yr 2013
\vol 32
\issue 19
\pages 92--100 

\Bibitem{Sol13}
\by Soldatov~A.~P.
\paper Integral representations in plane anisotropic elasticity
\jour AIP Conf. Proc. 
\vol 1570
\papernumber 276 
\yr 2013
\crossref {10.1063/1.4854766}

\RBibitem{Sol13}
\by Солдатов~А.~П. 
\paper Задача Дирихле для слабо связанных эллиптических систем на плоскости
\jour Диффер. уравн.
\yr 2013
\vol 49
\issue 6
\pages 734--745
\crossref{10.1134/S0374064113060058}
\transl
\paper Dirichlet problem for weakly coupled elliptic systems on the plane
\by Soldatov~A.~P.
\jour Differ. Equat.
\yr 2013
\vol 49
\issue 6
\pages 707--717
\crossref{10.1134/S0012266113060050}

\RBibitem{Sol13}
\by Жура~Н.~А., Солдатов~А.~П.
\crossref{10.7868/S0869565213340070}
\paper К решению обратной задачи Штурма–Лиувилля на всей оси
\jour Докл. Акад. наук
\yr 2013
\vol 453
\issue 4
\pages 368–372 
\transl
\paper To the solution of the inverse Sturm-Liouville problem on the real axis
\by \linebreak Zhura~N.~A., Soldatov~A.~P.
\jour Dokl. Math. 
\yr 2013
\vol 88
\issue 3
\pages 695--699
\crossref{10.1134/S1064562413060227}



\Bibitem{1}
\by Zhura~N.~A., Soldatov~A.~P.
\paper To inverse scattering problem of Gelfand, Levitan, Marchenko
\jour AIP Conf. Proc.
\vol 1570
\papernumber 298 
\yr 2013
\crossref{10.1063/1.4854769}


\Bibitem{1}
\by Soldatov~A.~P.
\paper Mixed problems for the Lavrent'ev--Bitsadze equation
\jour AIP Conf. Proc.
\vol 1497
\papernumber 199 
\yr 2012
\crossref{10.1063/1.4766786}


\RBibitem{ZhuSol12}
\by Жура~Н.~А., Солдатов~А.~П.
\paper \href{http://mi.mathnet.ru/smj2325}{Об асимптотике кусочно аналитической функции, удовлетворяющей контактным условиям}
\jour Сиб. матем. журн.
\yr 2012
\vol 53
\issue 5
\pages 1001--1006
\mathnet{http://mi.mathnet.ru/smj2325}
\mathscinet{http://www.ams.org/mathscinet-getitem?mr=3057682}
\elib{http://elibrary.ru/item.asp?id=17897049}
\transl
\paper On asymptotics of a piecewise analytic function satisfying contact conditions
\by \linebreak Zhura~N.~A.,  Soldatov~A.~P.
\jour Siberian Math. J.
\yr 2012
\vol 53
\issue 5
\pages 800--804
\crossref{10.1134/S0037446612050059}
\isi{http://gateway.isiknowledge.com/gateway/Gateway.cgi?GWVersion=2&SrcApp=PARTNER_APP&SrcAuth=LinksAMR&DestLinkType=FullRecord&DestApp=ALL_WOS&KeyUT=000310374900005}
\elib{http://elibrary.ru/item.asp?id=20498143}
\scopus{http://www.scopus.com/record/display.url?origin=inward&eid=2-s2.0-84868141053}

\RBibitem{Sol12}
\by Солдатов~А.~П.
\paper \href{http://mi.mathnet.ru/tm3402}{О задачах типа Дирихле для уравнения Лаврентьева--Бицадзе}
\inbook Дифференциальные уравнения и динамические системы
\bookinfo Сборник статей
\serial Тр. МИАН
\yr 2012
\vol 278
\pages 242--249
\publ МАИК «Наука/Интерпериодика»
\publaddr М.
\mathnet{http://mi.mathnet.ru/tm3402}
\mathscinet{http://www.ams.org/mathscinet-getitem?mr=3058799}
\transl
\paper On Dirichlet-type problems for the Lavrent'ev--Bitsadze equation
\by  Soldatov~A.~P.
\jour Proc. Steklov Inst. Math.
\yr 2012
\vol 278
\pages 233--240
\crossref{10.1134/S0081543812060223}
\isi{http://gateway.isiknowledge.com/gateway/Gateway.cgi?GWVersion=2&SrcApp=PARTNER_APP&SrcAuth=LinksAMR&DestLinkType=FullRecord&DestApp=ALL_WOS&KeyUT=000309861500022}
\scopus{http://www.scopus.com/record/display.url?origin=inward&eid=2-s2.0-84867372762}


\end{thebibliography}

\vfill


Редколлегия журнала <<Вестник Самарского государственного технического университета. Сер. Физико-математические науки>> поздравляет Александра Павловича с юбилеем и желает ему здоровья и дальнейших творческих успехов в научной и педагогической деятельности.

\vfill

\hfill {\it \rupid{38541}{А.~А.~Андреев}, \rupid{38375}{В.~П.~Радченко}, \rupid{62329}{Е.~А.~Козлова}}


\vfill


\AdditionalContent{Используемые материалы}{В заметке были использованы материалы  сайта \linebreak \href{https://www.bsu.edu.ru}{НИУ <<БелГУ>>} (\url{https://www.bsu.edu.ru/bsu/personal.php?ID=15120}) и персональной страницы А.~П.~Солдатова на портале \href{http://www.mathnet.ru}{Math-Net.Ru} (\url{http://www.mathnet.ru/person8905}). Фотография любезно предоставлена юбиляром.} 


\newpage

\selectlanguage{english}%
\normalsize 
\normalfont
\chead[\fancyplain{}{\selectlanguage{english}\theenauthor}]{\fancyplain{}{\selectlanguage{english} \theentitle}}
\rhead[]{\fancyplain{\EnJournal}{}}
\lhead[\fancyplain{\EnJournal}{}]{}
\cfoot[]{}
\lfoot[\thepage]{}
\rfoot[]{\thepage}
\renewcommand{\plainheadrulewidth}{0.7pt}
\renewcommand{\headrulewidth}{0.7pt}
\pagestyle{fancyplain}
\thispagestyle{plain}
\clear
\hskip-7mm
\begin{minipage}[t]{125mm}
\vskip.5mm
\noindent\themsc
\vskip2mm
\EnTitleInTitlepage
\end{minipage}
\par
\theenauthorsdetails
\begin{minipage}[t]{125mm}
\vskip4mm
{\renewcommand{\bfdefault}{bx}\noindent\small\textbf{\emph{\EnAuthorInTitlepage}}\renewcommand{\bfdefault}{b}}  
\vskip2mm
{\scriptsize \enaffilname}
\end{minipage}
\vskip2mm


\footnotesize

\smallskip

\selectlanguage{english}



The paper is devoted to Alexander Pavlovich Soldatov, the well known mathematician, and his contribution to the development of scientific theories and applications. 
Alexander Pavlovich has turned 70 years on the January, 18$^{\text{th}}$, so we give a short biographical background to his anniversary. 

Alexander Pavlovich Soldatov was a student of Novosibirsk State University, and graduated from with honours. 
He continued his education, defended a thesis in the famous V.~A.~Steklov Mathematical Institute of the USSR Academy of Sciences, became a Candidate of Physical and Mathematical Sciences and then Doctor habilitated of Physical and Mathematical Sciences (in Moscow State University). 

Today he is a professor, the Leading Researcher of Dorodnicyn Computing Centre  in Federal Research Center ``Computer Science and Control'' of RAS, and has a title of Honored Scientist of Russia. 
The studies of  Alexander Pavlovich in partial differential equations are well known to the specialists. 
He is an author of more then 170 scientific articles and four monographies. 
Here we give a list (not complete) of A.~P.~Soldatov publications for the last five years, that should help readers to see the modern interests of scientist.

Alexander Pavlovich Soldatov is an excellent teacher, also he works as reviewer for mathematical journals, develops the mathematical competitions for school and university students and organizes mathematical conferences including international.

In the paper we introduce the contribution of Alexander Pavlovich to the mathematics and education science, and give a short compilation of his remarkable scientific results.

\bigskip

\par\noindent
\theonlineengdate \par\noindent
\rule{125mm}{.7pt}\\

\footnotetext{
\vspace{-1em}\\
\noindent{\normalsize\bf\enArticleType}\smallskip\par\noindent
\OAlogo~\ccby\ The content is published under the terms of the \href{http://creativecommons.org/licenses/by/4.0/}{Creative Commons \mbox{Attribution} 4.0 International} License (\url{http://creativecommons.org/licenses/by/4.0/})\selectlanguage{english}
\vspace{.3em}\\           
\EnCitation}




\selectlanguage{russian}

\newpage

%\end{document}
%
